\documentclass[11pt,a4paper]{article}

\usepackage[utf8]{inputenc}

\usepackage[T1]{fontenc}

\usepackage[polish]{babel}

\usepackage{lmodern}

\usepackage{geometry}

\geometry{a4paper,margin=27mm}

\usepackage{microtype}

\usepackage[hidelinks]{hyperref}

% Portfolio: MSC-AR-P05; source: prospective ARC-XR research scope.

\title{Weryfikacja poprawności sekwencji manualnej w AR na podstawie dowodów wizualnych}

\author{}

\date{Wersja 0.1.0, 21 listopada 2026}

\begin{document}

\maketitle

\begin{abstract}

Praca dotyczy zagadnienia: Weryfikacja poprawności sekwencji manualnej w AR na podstawie dowodów wizualnych. Czy prosta pamięć sekwencji zwiększa wykrywalność pominiętych kroków względem analizy pojedynczej klatki? Proponowane porównanie obejmuje porównać VLM oceniający każdą obserwację niezależnie z wariantem uwzględniającym bieżący stan prostej procedury. Oczekiwany rodzaj wyniku stanowi czułość wykrycia pominiętego kroku przy ustalonym limicie fałszywych alarmów na sekwencję. Tekst jest wstępnym szkic manuskryptu; nie przedstawia niezweryfikowanych danych ani wniosków z badań.

\end{abstract}

\section{Wprowadzenie}

Rozwój adaptacyjnych systemów XR stwarza możliwość badania interakcji człowieka z informacją przestrzenną, ale samo zastosowanie AR, VR lub modeli sztucznej inteligencji nie przesądza o przewadze poznawczej czy ergonomicznej. W rozpatrywanym zagadnieniu, Weryfikacja poprawności sekwencji manualnej w AR na podstawie dowodów wizualnych, kluczowe jest rozdzielenie właściwości systemu od zmienności zachowania użytkownika. Porównanie należy zaprojektować tak, aby mierzyć jednoznaczne zjawisko, a nie tylko ogólną ocenę technologii.

\section{Cel pracy i pytanie badawcze}

Czy prosta pamięć sekwencji zwiększa wykrywalność pominiętych kroków względem analizy pojedynczej klatki? Najpierw oceniana jest detekcja błędu w benchmarku offline; test użyteczności przestrzennych instrukcji można przeprowadzić jako osobne badanie.

\section{Materiał i metody}

Porównać VLM oceniający każdą obserwację niezależnie z wariantem uwzględniającym bieżący stan prostej procedury. Wymagane jest zdefiniowanie jednostki analizy, zasad doboru przypadków, wersji scenariusza i źródeł pomiaru. Rozkład warunków oraz reguły wykluczenia prób należy określić przed pomiarem. W badaniach z udziałem ludzi niezbędna jest uprzednia właściwa ocena etyczna i świadoma zgoda. W testach wyłącznie symulacyjnych należy oddzielić prawdę symulacyjną od obserwacji empirycznej.

\subsection{Wynik główny i kryteria jakości}

Czułość wykrycia pominiętego kroku przy ustalonym limicie fałszywych alarmów na sekwencję. Pomiar powinien zachowywać jawne jednostki, znaczniki czasu i reguły oceny poprawności. Przy modelach predykcyjnych stosować rozłączne partycje oceny; przy eksperymentach interwencyjnych oddzielić skuteczność predykcji od przyczynowego wpływu interfejsu.

\section{Plan analizy}

Analiza zdarzeniowa dla niezależnych procedur i sesji, z przedziałami ufności dla czułości, alarmów i opóźnienia. W wynikach należy raportować wielkość efektu lub błąd estymacji z przedziałem niepewności, liczbą obserwacji i odsetkiem odrzuconych prób. Analizy dodatkowe wymagają jasnego oznaczenia jako eksploracyjne. Nie wybierać progów decyzyjnych ani modelu na podstawie obserwacji ze zbioru testowego.

\section{Wyniki}

W obecnej wersji nie raportuje się wyników pomiarów, tabel statystycznych ani parametrów wyznaczonych z nieudokumentowanych danych. Nie stwierdza się w tym opracowaniu, że eksperyment został wykonany lub zgłoszony do czasopisma.

\section{Dyskusja}

Interpretacja przyszłych wyników musi uwzględnić alternatywne mechanizmy i zależność wniosków od jakości danych, warunków zadania oraz ograniczeń sprzętu. Wartość potencjalnego rozwiązania zależy nie tylko od trafności modelu, lecz także od jego opóźnienia, stabilności i wpływu na zachowanie człowieka. Wyniki testów syntetycznych nie mogą być automatycznie uogólniane na populację uczestników.

\section{Ograniczenia badania}

Błąd anotacji granic kroków i częściowa widoczność dłoni mogą powodować pozorną poprawę wykrywania. Weryfikacja etapów procedury, nie grounding obiektów P04 ani faza podania wskazówki MSC-AR-P25. Żadna część szkicu nie dowodzi przeprowadzenia zbierania danych, uzyskania zgód etycznych ani gotowości manuskryptu do zgłoszenia.

\section{Wnioski}

Przedstawiono problem, kontrast i sposób jego planowanej oceny. Rozstrzygnięcie postawionego pytania wymaga przygotowania protokołu, zweryfikowanych źródeł oraz poprawnie udokumentowanych wyników. Bez tych materiałów nie można sformułować konkluzji empirycznych.

\end{document}

